Benchmark contamination in Large Language Models (LLMs) undermines evaluation validity, yet most detection methods require privileged access to training data, model internals, or output probabilities.
We present the first controlled characterization of Item Response Theory (IRT) person-fit statistics for detecting supervised fine-tuning (SFT) contamination, requiring only binary response patterns and calibrated four-parameter logistic item parameters; diagnostic $\theta$-anchoring requires a clean baseline.
A contaminated LLM produces aberrant response patterns analogous to a test-taker with item preknowledge.
Using $\ell_z^*$ in a matched-difference design ($\Delta\ell_z^*$) across three 7--8B model families ($n{=}10$ clean seeds per family) on MMLU (14,042 items) and ARC-Challenge (295 items), we identify three boundaries:
(1) model-specific $\theta$-shift dynamics, where contamination-induced ability inflation absorbs the person-fit signal; $\theta$-anchoring amplifies detection by 2.4$\times$ while reducing coefficient of variation from 60\% to 19\%;
(2) a detection-diagnosis tradeoff where simpler statistics yield stronger detection but lack $\ell_z^*$'s diagnostic capabilities (seen/unseen decomposition, mechanistic $\theta$-shift analysis);
and (3, exploratory) ability-dependent non-monotonic dose-response under free-$\theta$ estimation, driven by $\theta$-shift absorption; $\theta$-anchored analysis reveals stable positive dose-response.
These findings establish when and why binary response patterns carry contamination information beyond aggregate accuracy.
